%!TEX TS-program = lualatex
%!TEX encoding = UTF-8 Unicode

\documentclass[t]{beamer}

%%%% HANDOUTS For online Uncomment the following four lines for handout
%\documentclass[t,handout]{beamer}  %Use this for handouts.
%\includeonlylecture{student}
%\usepackage{handoutWithNotes}
%\pgfpagesuselayout{3 on 1 with notes}[letterpaper,border shrink=5mm]


%%% Including only some slides for students.
%%% Uncomment the following line. For the slides,
%%% use the labels shown below the command.

%% For students, use \lecture{student}{student}
%% For mine, use \lecture{instructor}{instructor}


% FONTS
\usepackage{fontspec}
\def\mainfont{Linux Biolinum O}
\setmainfont[Ligatures={Common,TeX}, Contextuals={NoAlternate}, BoldFont={* Bold}, ItalicFont={* Italic}, Numbers={OldStyle}]{\mainfont}
\setsansfont[Ligatures={Common,TeX}, Scale=MatchLowercase, Numbers=OldStyle]{Linux Biolinum O} 
\usepackage{microtype}

\usepackage{graphicx}
	\graphicspath{{/Users/mtaylor/pictures/teach/163/lecture/}
	{/Users/mtaylor/pictures/teach/common/}} % set of paths to search for images

\usepackage{multicol}
\usepackage{booktabs}
%\usepackage{textcomp}

\usepackage{tikz}
	\tikzstyle{every picture}+=[remember picture,overlay]
	\usetikzlibrary{trees}
%	\usetikzlibrary{arrows}

\mode<presentation>
{
  \usetheme{Lecture}
  \setbeamercovered{invisible}
%  \setbeamertemplate{items}[square]
}


\begin{document}

\begin{frame}{Our goals for this lecture are to learn}
	
	\hangpara what is \highlight{descent with modification,}

	\hangpara	the process of \highlight{natural selection,} and

	\hangpara what is meant by \highlight{relative fitness.}

\end{frame}

%
{
\usebackgroundtemplate{\includegraphics[width=\paperwidth]{hms_beagle}}
\begin{frame}[b]

\hfill\tiny\textcolor{white}{Wikimedia, public domain.}
\end{frame}
}
%
{
\usebackgroundtemplate{\includegraphics[width=\paperwidth]{voyage_beagle}}
\begin{frame}[t]{\textcolor{white}{The \textsc{h.m.s.} Beagle sailed around the world during 1831--1836.}}
\end{frame}
}
%
{
\usebackgroundtemplate{\includegraphics[width=\paperwidth]{giant_sloth}}
\begin{frame}[t]{In South America, Darwin witnessed\ldots}

	\hangpara fossils of a giant sloth and other animals,

	\hangpara marine shells on a mountain top, and

	\hangpara a magnitude 8 earthquake,

	\hangpara the Galapagos Islands.

	\vskip0pt plus 1filll

	\hfill\tiny ДиБгд, Wikimedia, public domain.
\end{frame}
}
%
{
\usebackgroundtemplate{\includegraphics[width=\paperwidth]{darwins_finches}}
\begin{frame}[b]

	\hfill\tiny Wikimedia, public domain.
\end{frame}
}
%
{
\usebackgroundtemplate{\includegraphics[width=\paperwidth]{darwin_descent_quote.jpg}}
\begin{frame}[plain]
\end{frame}
}
%
{
\usebackgroundtemplate{\includegraphics[width=\paperwidth]{darwin_phylogeny}}
\begin{frame}[t]
\end{frame}
}
%
{
\usebackgroundtemplate{\includegraphics[width=\paperwidth]{artificial_selection}}
\begin{frame}[t]{}
\end{frame}
}
%

{
\usebackgroundtemplate{\includegraphics[width=\paperwidth]{origin_of_species_cover}}
\begin{frame}[t]

	\vspace*{3\baselineskip}

	\hangpara\parbox[t]{2.25in}{\raggedright %
There is grandeur in this view of life, with its several powers, having been originally breathed into a few forms or into one; and that, whilst this planet has gone cycling on according to the fixed law of gravity, from so simple a beginning endless forms most beautiful and most wonderful have been, and are being, evolved.}

	\vskip0pt plus 1filll
	
	\hfill\tiny Wikimedia, public domain.
\end{frame}
}
%
{
\usebackgroundtemplate{\includegraphics[width=\paperwidth]{origin_of_species_cover}}
\begin{frame}[t]

	\vspace*{3\baselineskip}
	
	\hangpara Darwin's book had two ideas:
	
	\hangpara\hspace*{1em}\highlight{descent with modification} 
	
	\hangpara\hspace*{1em}\highlight{natural selection.}

	\vskip0pt plus 1filll
	
	\hfill\tiny Wikimedia, public domain.
\end{frame}
}
%
\begin{frame}{\highlight{Descent with modification} means that}
\vspace{2ex}
\hspace{1em}all organisms \highlight{descend} from a common ancestor, and
\pause

\vspace{1ex}
\hspace{1em}the descendants are \highlight{modified} from the ancestral form.

\vspace{1ex}
\pause

\centering
\begin{tikzpicture}
%	[every node/.style={orange6, rounded corners},
	[every node/.append style={rounded corners}]
	\node at (0,-3.5) [color=orange7] {\textbf{ancestor}}
		[edge from parent fork up, grow=up,
		edge from parent/.style={very thick, draw}]
		child [sibling distance=20mm] {
			child {node [color=green6] {\textbf{desc. \textsc{c}}}}
			child {node [color=blue7] {\textbf{desc. \textsc{b}}}}
		}
		child [sibling distance=40mm] {
			child {node [color=orange6] {\textbf{desc. \textsc{a} }}}
			%child {node {D}}
		};
\end{tikzpicture}

\end{frame}
%
\lecture{instructor}{instructor}
{
\usebackgroundtemplate{\includegraphics[width=\paperwidth]{descent_with_modification1}}
\begin{frame}[b]

\tiny\textcopyright Pearson Education, Inc.
\end{frame}
}
%
\lecture{student}{student}
{
\usebackgroundtemplate{\includegraphics[width=\paperwidth]{descent_with_modification2}}
\begin{frame}[b]

\tiny\textcopyright Pearson Education, Inc.
\end{frame}
}
%
%
{
\usebackgroundtemplate{\includegraphics[width=\paperwidth]{darwin_natural_selection_quote_new}}
\begin{frame}
\end{frame}
}
%
\lecture{instructor}{instructor}
{
\usebackgroundtemplate{\includegraphics[width=\paperwidth]{darwin_natural_selection_quote_highlight}
}
\begin{frame}[plain]
\end{frame}
}

\lecture{student}{student}
{
\usebackgroundtemplate{\includegraphics[width=\paperwidth]{individuals_vary}
}
\begin{frame}[b,plain]
\hfill\textcolor{white}{\tiny \textcopyright Entomart, Wikimedia}
\end{frame}
}
%
{
\usebackgroundtemplate{\includegraphics[width=\paperwidth]{phenotype_environment}
}
\begin{frame}[b,plain]{\textit{Only} genetic variation is subject to selection.}
\hfill\tiny Brakefield and Finch 1999. BioEssays 21:391.
\end{frame}
}
%
{
\usebackgroundtemplate{\includegraphics[width=\paperwidth]{overproduce_offspring}}
\begin{frame}[b]
\tiny\textcolor{white}{Silke Baron, Wikimedia \ccby{2}}
\end{frame}
}
%
{
\usebackgroundtemplate{\includegraphics[width=\paperwidth]{not_all_will_survive}}
\begin{frame}[b]
\hfill\textcolor{white}{\tiny Pierre Dalous, Wikimedia, \ccbysa{}}
\end{frame}
}
%
{
\usebackgroundtemplate{\includegraphics[width=\paperwidth]{traits_determine_survival}}
\begin{frame}[b]
\hfill\textcolor{white}{\tiny \textcopyright Dustin Siegel, with permission.}
\end{frame}
}
%
{
\usebackgroundtemplate{\includegraphics[width=\paperwidth]{more_offspring_baby_spiders}}
\begin{frame}[b]
\hfill\textcolor{white}{\tiny C. Frank Starmer, \href{http://frank.itlab.us}{frank.itlab.us},\ccbyncsa{}}
\end{frame}
}
%
{
\usebackgroundtemplate{\includegraphics[width=\paperwidth]{better_traits_dandelion}}
\begin{frame}[b]
\hfill\textcolor{white}{\tiny David Carillet, \href{http://www.wildretina.com}{www.wildretina.com}, \ccbyncsa{3}}
\end{frame}
}


% Cycle of Natural Selection
\begin{frame}[c]
	\begin{center}
	
	\vspace{1.5cm}
	\begin{tikzpicture}
			[scale=0.6, transform shape,
			myNode/.style={circle, inner sep=2mm, minimum size=4cm, draw=black, fill=blue3},
			myLine/.style={line width=2pt, draw=black!70!white, ->},
			textNode/.style={text width=3.5cm, align=flush center, font=\Large}]
		\node (adapt) at (0,0) {\textcolor{orange6}{\textbf{\Huge Adaptation}}};
		\foreach \r in {60,120,...,360} {%
			\begin{scope}[shift={(\r:6)}]
				\node[myNode] (n\r) {};
			\end{scope};
		};
		\draw [myLine] (n60) to (n360);
		\draw [myLine] (n360) to (n300);
		\draw [myLine] (n300) to (n240);
		\draw [myLine] (n240) to (n180);
		\draw [myLine] (n180) to (n120);
		\draw [myLine] (n120) to (n60);
	
		\node [textNode] at (n60) {Individuals \highlight{vary} in populations};
		\node [textNode] at (n360) {Populations \highlight{overproduce} offspring};
		\node [textNode] at (n300) {Not all \highlight{survive}};
		\node [textNode] at (n240) {\highlight{Traits} determine survival};
		\node [textNode] at (n180) {More survival, \highlight{more offspring}};
		\node [textNode] at (n120) {Better traits \highlight{spread} through population};
	
	\end{tikzpicture}
	
	\end{center}
\end{frame}

\setbeamercovered{invisible}
\setbeamercovered{again covered={\opaqueness<1->{15}}}
\begin{frame}[t]{Natural selection leads to descent with modification.}


	\hangpara \textcolor<2-4>{gray}{Natural selection causes variation in \only<1,4>{\highlight}{reproductive success} among individuals that vary in  \only<1>{\highlight}{heritable traits.}}

	\pause
	
	\hangpara	 \textcolor<3->{gray}{Traits that increase reproductive success are called  \only<2>{\highlight}{adaptations.}}
	
	\pause 
	
	\hangpara	 \textcolor<4->{gray}{Adaptations are traits that have been  \only<3>{\highlight}{modified} from the ancestral form.}

	\pause 
		
	\vspace{3ex}

	\hangpara Let us study variation in reproductive success more closely.


\end{frame}

\setbeamercovered{invisible}
\begin{frame}[t]{True or false?}

	\hangpara \textcolor<2->{gray}{Darwin coined the phrase “survival of the fittest.”}
	
	\pause
	
	\vspace*{2ex}
	
	\hangpara\alt<handout>{}{\highlight{False.} The phrase was coined by Herbert Spencer.}
	
	\hangpara\alt<handout>{}{Darwin did not use the phrase until the 5th edition of \textit{On the Origin of Species.}}

\end{frame}

\begin{frame}[t]{“Survival of the fittest” is a circular argument.}

	\vspace*{2ex}
	
	\begin{center}
	\begin{tikzpicture}
		[textNode/.style={text width=2.75cm, align=flush center, font=\huge},
		smallNode/.style={text width=2cm, font=\large},
		->, ultra thick]
	
		\node [textNode] (survival) at (0,0) {Survival};

		\node [textNode] at (0,-1.5) {\small\textcolor{gray!70!black}{of the}};

		\node [textNode] (fittest) at (0,-3) {Fittest};
		
		\draw (survival.east) [bend left=60]  to (fittest.east);

		\draw (fittest.west) [bend left=60]  to (survival.west);

		\pause
		
		\node [smallNode, align=flush left] at (3.5,-1.5) {Which survive?};

		\pause
		
		\node [smallNode, align=flush right] at (-3.5,-1.5) {Which are the fittest?};
		
		\pause
		
		\node [align=left] at (-1.5,-5) {Fitness must be measured independently of survival.};

	\end{tikzpicture}
	\end{center}
	
%	\pause
%	
%	\vspace*{20ex}
%	
%	\hangpara We need an independent way to measure fitness.
	
\end{frame}


%% REPLACE THESE THREE WITH APPROPRIATE images
{
\usebackgroundtemplate{\includegraphics[width=\paperwidth]{fitness_offspring}}
\begin{frame}[b]
\hfill\textcolor{white}{\tiny Guido Bohne, Flickr \ccbysa{2}}
\end{frame}
}
%
{
\usebackgroundtemplate{\includegraphics[width=\paperwidth]{fitness_is_relative}}
\begin{frame}[b]
\textcolor{white}{\tiny Michael S. Taylor, \ccbysa{4}}
\end{frame}
}
%
{
\usebackgroundtemplate{\includegraphics[width=\paperwidth]{fitness_not_physical}}
\begin{frame}[b]
\hfill\textcolor{white}{\tiny Ian Carroll, Flickr, \ccby{2}}
\end{frame}
}
%
{
\usebackgroundtemplate{\includegraphics[width=\paperwidth]{savings_account}}
\begin{frame}[b]
\hfill\tiny 401(k) 2012, Flickr, \ccbysa{2}
\end{frame}
}
%
{
\usebackgroundtemplate{\includegraphics[width=\paperwidth]{financial_analogy}}
\begin{frame}[b]
%\hfill\textcolor{white}{\tiny Guido Bohne, \ccbysa{2}}
\end{frame}
}
%

%\includegraphics[page=\arabic{4},width=.95\textwidth]{a.pdf}

\lecture{student}{student}
%
\begin{frame}[t]{Our goals for this lecture are to learn}

	
	\hangpara	the difference between \highlight{phenotype} and \highlight{genotype,} and

	\hangpara how phenotype changes in response to different \highlight{modes of selection.}
	
\end{frame}

%
%{
%\usebackgroundtemplate{\includegraphics[width=\paperwidth]{individuals_vary}}
%\begin{frame}[b]
%\hfill\textcolor{white}{\tiny \copyright\, Entomart, Wikimedia}
%\end{frame}
%}
% Why do populations vary?
%
\lecture{student}{student}

\begin{frame}[t]{If natural selection changes the phenotype, what also had to change?}

\hangpara Three \highlight{modes of selection} describe \emph{how} the phenotype changes in a population.

\hangpara The modes do not say \emph{why} the phenotype changes.

\end{frame}
%
\begin{frame}{\highlight{Directional} selection shifts the population toward \emph{one} extreme phenotype.}
	\centering
	\begin{columns}[T]
		\begin{column}{0.5\textwidth}
			\centering
			\includegraphics[width=0.9\textwidth]{mode_original}
		\end{column}
		\begin{column}{0.5\textwidth}
			\pause\includegraphics[width=0.9\textwidth]{mode_directional}
		\end{column}
	\end{columns}
\end{frame}
%
{
\usebackgroundtemplate{\includegraphics[width=\paperwidth]{scissortail_flycatcher} }
\begin{frame}[b]{Females prefer longer tails in male scissortail flycatchers.}

	\hfill \tiny TexasEagle, Flickr, \ccbync{2}

\end{frame}
}
%
\begin{frame}{\highlight{Disruptive} selection shifts the population toward \emph{both} extreme phenotypes.}
	\centering
	\begin{columns}[T]
		\begin{column}{0.5\textwidth}
			\centering
			\includegraphics[width=0.9\textwidth]{mode_original}
		\end{column}
		\begin{column}{0.5\textwidth}
			\includegraphics[width=0.9\textwidth]{mode_disruptive}
		\end{column}
	\end{columns}
\end{frame}
%
{
\usebackgroundtemplate{\includegraphics[width=\paperwidth]{geospiza_fortis} }
\begin{frame}[b]{Beak size is diverging in \textit{Geospiza fortis.}}

	\tiny\textcolor{white}{Mark Putney, Flickr, \ccbysa{2}} \pause\hfill\includegraphics[width=0.45\textwidth]{geospiza_fortis_disruptive}

\end{frame}
}
%
\begin{frame}{\highlight{Stabilizing} selection shifts the population toward the intermediate phenotype.}
	\centering
	\begin{columns}[T]
		\begin{column}{0.5\textwidth}
			\centering
			\includegraphics[width=0.9\textwidth]{mode_original}
		\end{column}
		\begin{column}{0.5\textwidth}
			\includegraphics[width=0.9\textwidth]{mode_stabilizing}
		\end{column}
	\end{columns}
\end{frame}
%
{
\usebackgroundtemplate{\includegraphics[width=\paperwidth]{balancing_selection} }
\begin{frame}[b]{}

\tiny \copyright\,Pearson Education, Inc.
\end{frame}
}

\end{document}
